The Hilbert calculus \(\K\) consists of all uniform substitution
instances of classically valid propositional formulas, modus ponens, the
distribution schema
\[
  \Box(A\to B)\to(\Box A\to\Box B),
  \tag{K}
\]
and necessitation
\[
  \frac{\derives_{\K}A}{\derives_{\K}\Box A}\;\textsc{Nec}.
\]
Necessitation applies to theorems, not to a formula depending on an
undischarged premise.  Define
\[
\begin{array}{rcl}
  \mathsf T&=&\K+(\Box A\to A),\\
  \Sfour&=&\mathsf T+(\Box A\to\Box\Box A),\\
  \Sfive&=&\mathsf T+(\Diamond A\to\Box\Diamond A).
\end{array}
\]
For \(\mathsf L\in\{\K,\mathsf T,\Sfour,\Sfive\}\) and a premise set
\(\Gamma\), an \(\mathsf L\)-derivation from \(\Gamma\) is a finite sequence
whose lines are members of \(\Gamma\), instances of the axiom schemas of
\(\mathsf L\), consequences of earlier lines by modus ponens, or formulas
\(\Box A\) for which \(A\) has a separate no-premise
\(\mathsf L\)-derivation.  Write \(\Gamma\derives_{\mathsf L}A\) when such a
sequence ends in \(A\).  Thus necessitation remains theorem-only even in a
derivation from premises.  This definition supplies the premise-sensitive
standing needed by strong completeness.

